\section*{Capitulo 10. Estadistica, probabilidad e inferencia}
\addcontentsline{toc}{section}{Capitulo 10. Estadistica, probabilidad e inferencia}

\begin{shortanswer}{[GX-C10.S01-01]}
\(P(U\mid S)=\dfrac{18}{26}=\dfrac9{13}\).
\end{shortanswer}

\begin{shortanswer}{C10.S01 practica autonoma}
\begin{enumerate}[label=\textbf{PX-C10.S01-\arabic*:}, leftmargin=*]
  \item \(18/40=0.45\)
  \item Supera: \(26/40=0.65\); no supera: \(14/40=0.35\)
  \item La tabla muestra asociacion, pero no controla posibles variables de confusion
\end{enumerate}
\end{shortanswer}

\begin{shortanswer}{[GX-C10.S02-01]}
Si. Una correlacion lineal casi nula no descarta una relacion no lineal, como una forma en U.
\end{shortanswer}

\begin{shortanswer}{C10.S02 practica autonoma}
\begin{enumerate}[label=\textbf{PX-C10.S02-\arabic*:}, leftmargin=*]
  \item Asociacion lineal negativa fuerte
  \item No; una correlacion perfecta no demuestra por si sola causalidad
  \item Puede aumentar, reducir o incluso cambiar el signo de \(r\)
\end{enumerate}
\end{shortanswer}

\begin{shortanswer}{[GX-C10.S03-01]}
El modelo aumenta \(0.8\) gramos por cada hora adicional.
\end{shortanswer}

\begin{shortanswer}{C10.S03 practica autonoma}
\begin{enumerate}[label=\textbf{PX-C10.S03-\arabic*:}, leftmargin=*]
  \item \(18\)
  \item \(0.8\)
  \item La relacion observada puede cambiar fuera del intervalo de datos
\end{enumerate}
\end{shortanswer}

\begin{shortanswer}{[GX-C10.S04-01]}
Muestreo aleatorio estratificado por curso y grupo, con seleccion aleatoria dentro de cada
estrato.
\end{shortanswer}

\begin{shortanswer}{C10.S04 practica autonoma}
\begin{enumerate}[label=\textbf{PX-C10.S04-\arabic*:}, leftmargin=*]
  \item Poblacion: \(900\) estudiantes; muestra: \(120\)
  \item Excluye o infrarrepresenta a quienes no frecuentan la biblioteca
  \item En la muestra, el \(60\%\) apoya la medida; generalizar exige justificar el muestreo
\end{enumerate}
\end{shortanswer}

\begin{shortanswer}{[GX-C10.S05-01]}
\(137/500=0.274\).
\end{shortanswer}

\begin{shortanswer}{C10.S05 practica autonoma}
\begin{enumerate}[label=\textbf{PX-C10.S05-\arabic*:}, leftmargin=*]
  \item \(0.63\)
  \item \(0.7\)
  \item \(\overline A\cap\overline B=\overline{A\cup B}\)
\end{enumerate}
\end{shortanswer}

\begin{shortanswer}{[GX-C10.S06-01]}
\(5\cdot4\cdot3=60\) codigos.
\end{shortanswer}

\begin{shortanswer}{C10.S06 practica autonoma}
\begin{enumerate}[label=\textbf{PX-C10.S06-\arabic*:}, leftmargin=*]
  \item \(\binom{10}{2}=45\)
  \item \(4!=24\)
  \item \(6/36=1/6\)
\end{enumerate}
\end{shortanswer}

\begin{shortanswer}{[GX-C10.S07-01]}
Si, porque \(P(A\cap B)=0.2=P(A)P(B)\).
\end{shortanswer}

\begin{shortanswer}{C10.S07 practica autonoma}
\begin{enumerate}[label=\textbf{PX-C10.S07-\arabic*:}, leftmargin=*]
  \item \(0.18/0.6=0.3\)
  \item \(0.7\cdot0.02+0.3\cdot0.05=0.029=2.9\%\)
  \item Un positivo tambien puede proceder de un falso positivo; debe considerarse la
  prevalencia mediante Bayes
\end{enumerate}
\end{shortanswer}

\begin{shortanswer}{[CH-C10.S07-01]}
La entrega debe incluir datos y unidades trazables, nube de puntos, modelos lineal y cuadratico,
\(R^2\), residuos, una prediccion justificada y limitaciones sobre extrapolacion y causalidad.
\end{shortanswer}
